\documentclass{amsart}
% For dealing with references we use the comment environment
\usepackage{verbatim}
\newenvironment{reference}{\comment}{\endcomment}
%\newenvironment{reference}{}{}
\newenvironment{slogan}{\comment}{\endcomment}
\newenvironment{history}{\comment}{\endcomment}

% For commutative diagrams we use Xy-pic
\usepackage[all]{xy}

% We use 2cell for 2-commutative diagrams.
\xyoption{2cell}
\UseAllTwocells

% We use multicol for the list of chapters between chapters
\usepackage{multicol}

% This is generally recommended for better output
\usepackage{lmodern}
\usepackage[T1]{fontenc}

% For cross-file-references

% Package for hypertext links:
\usepackage{hyperref}

% For any local file, say "hello.tex" you want to link to please

% Theorem environments.
%
\theoremstyle{plain}
\newtheorem{theorem}[subsection]{Theorem}
\newtheorem{proposition}[subsection]{Proposition}
\newtheorem{lemma}[subsection]{Lemma}

\theoremstyle{definition}
\newtheorem{definition}[subsection]{Definition}
\newtheorem{example}[subsection]{Example}
\newtheorem{exercise}[subsection]{Exercise}
\newtheorem{situation}[subsection]{Situation}

\theoremstyle{remark}
\newtheorem{remark}[subsection]{Remark}
\newtheorem{remarks}[subsection]{Remarks}

\numberwithin{equation}{subsection}

% Macros
%
\def\lim{\mathop{\mathrm{lim}}\nolimits}
\def\colim{\mathop{\mathrm{colim}}\nolimits}
\def\Spec{\mathop{\mathrm{Spec}}}
\def\Hom{\mathop{\mathrm{Hom}}\nolimits}
\def\Ext{\mathop{\mathrm{Ext}}\nolimits}
\def\SheafHom{\mathop{\mathcal{H}\!\mathit{om}}\nolimits}
\def\SheafExt{\mathop{\mathcal{E}\!\mathit{xt}}\nolimits}
\def\Sch{\mathit{Sch}}
\def\Mor{\mathop{\mathrm{Mor}}\nolimits}
\def\Ob{\mathop{\mathrm{Ob}}\nolimits}
\def\Sh{\mathop{\mathit{Sh}}\nolimits}
\def\NL{\mathop{N\!L}\nolimits}
\def\CH{\mathop{\mathrm{CH}}\nolimits}
\def\proetale{{pro\text{-}\acute{e}tale}}
\def\etale{{\acute{e}tale}}
\def\QCoh{\mathit{QCoh}}
\def\Ker{\mathop{\mathrm{Ker}}}
\def\Im{\mathop{\mathrm{Im}}}
\def\Coker{\mathop{\mathrm{Coker}}}
\def\Coim{\mathop{\mathrm{Coim}}}

% Boxtimes
%
\DeclareMathSymbol{\boxtimes}{\mathbin}{AMSa}{"02}

%
% Macros for moduli stacks/spaces
%
\def\QCohstack{\mathcal{QC}\!\mathit{oh}}
\def\Cohstack{\mathcal{C}\!\mathit{oh}}
\def\Spacesstack{\mathcal{S}\!\mathit{paces}}
\def\Quotfunctor{\mathrm{Quot}}
\def\Hilbfunctor{\mathrm{Hilb}}
\def\Curvesstack{\mathcal{C}\!\mathit{urves}}
\def\Polarizedstack{\mathcal{P}\!\mathit{olarized}}
\def\Complexesstack{\mathcal{C}\!\mathit{omplexes}}
% \Pic is the operator that assigns to X its picard group, usage \Pic(X)
% \Picardstack_{X/B} denotes the Picard stack of X over B
% \Picardfunctor_{X/B} denotes the Picard functor of X over B
\def\Pic{\mathop{\mathrm{Pic}}\nolimits}
\def\Picardstack{\mathcal{P}\!\mathit{ic}}
\def\Picardfunctor{\mathrm{Pic}}
\def\Deformationcategory{\mathcal{D}\!\mathit{ef}}

\setlength{\emergencystretch}{3em}
\begin{document}
\title[Stacks Project, Tag 0A3G]{641 --- Cohomological dimension bounds for spectral spaces}
\maketitle
\noindent Copyright (C) 2005--2025 Johan de Jong. Stacks Project authors. Source: \href{https://stacks.math.columbia.edu/tag/0A3G}{Tag 0A3G}.

\noindent Editorial difficulty note (not author text). Difficulty H3: The proof first constructs top-degree extensions using neighborhoods of specialization loci and Mayer-Vietoris, then uses that surjectivity to force an injective-resolution kernel to be acyclic. The simultaneous dimensional induction, Cech vanishing, and support bounds are specialized research-level topology..

\noindent Editorial provenance note (not author text). History: excerpt from revision \texttt{a0444\allowbreak 6e57e\allowbreak c1fbc\allowbreak 252a8\allowbreak 71afc\allowbreak ec775\allowbreak 2fb28\allowbreak 07b14}, cohomology.tex, lines 3939--4057. Reference presentation was adapted; original-author context and core support blocks were retained or added. The mathematical wording is unchanged. Licensed under GNU FDL 1.2 or later; no invariant sections or cover texts. See ../licenses/COPYING. Context blocks reproduce author definitions and supporting results; they are not additional counted problems.

\begin{lemma}
\label{lemma-vanishing-ravi}
Let $X$ be a topological space. Let $\mathcal{F}$ be an abelian sheaf
on $X$. Let $\mathcal{U} : U = \bigcup_{i \in I} U_i$ be an
open covering. Assume the restriction mappings
$\mathcal{F}(U) \to \mathcal{F}(U')$ are surjective
for $U'$ an arbitrary union of opens of the form $U_{i_0 \ldots i_p}$.
Then $\check{H}^p(\mathcal{U}, \mathcal{F})$
vanishes for $p > 0$.
\end{lemma}

\begin{proof}
Let $Y$ be the set of nonempty subsets of $I$. We will use the letters
$A, B, C, \ldots$ to denote elements of $Y$, i.e., nonempty subsets of $I$.
For a finite nonempty subset $J \subset I$ let
$$
V_J = \{A \in Y \mid J \subset A\}
$$
This means that $V_{\{i\}} = \{A \in Y \mid i \in A\}$ and
$V_J = \bigcap_{j \in J} V_{\{j\}}$.
Then $V_J \subset V_K$ if and only if $J \supset K$.
There is a unique topology on $Y$ such that the collection of
subsets $V_J$ is a basis for the topology on $Y$. Any open is of the form
$$
V = \bigcup\nolimits_{t \in T} V_{J_t}
$$
for some family of finite subsets $J_t$. If $J_t \subset J_{t'}$
then we may remove $J_{t'}$ from the family without changing $V$.
Thus we may assume there are no inclusions among the $J_t$.
In this case the minimal elements of $V$ are the sets $A = J_t$.
Hence we can read off the family $(J_t)_{t \in T}$ from the open $V$.

\medskip\noindent
We can completely understand open coverings in $Y$. First, because
the elements $A \in Y$ are nonempty subsets of $I$ we have
$$
Y = \bigcup\nolimits_{i \in I} V_{\{i\}}
$$
To understand other coverings, let $V$ be as above and let $V_s \subset Y$
be an open corresponding to the family $(J_{s, t})_{t \in T_s}$. Then
$$
V = \bigcup\nolimits_{s \in S} V_s
$$
if and only if for each $t \in T$ there exists an $s \in S$ and
$t_s \in T_s$ such that $J_t = J_{s, t_s}$. Namely, as the family
$(J_t)_{t \in T}$ is minimal, the minimal element $A = J_t$
has to be in $V_s$ for some $s$, hence $A \in V_{J_{t_s}}$ for some
$t_s \in T_s$. But since $A$ is also minimal in $V_s$ we conclude
that $J_{t_s} = J_t$.

\medskip\noindent
Next we map the set of opens of $Y$ to opens of $X$. Namely, we send
$Y$ to $U$, we use the rule
$$
V_J \mapsto U_J = \bigcap\nolimits_{i \in J} U_i
$$
on the opens $V_J$, and we extend it to arbitrary opens $V$ by the rule
$$
V = \bigcup\nolimits_{t \in T} V_{J_t}
\mapsto
\bigcup\nolimits_{t \in T} U_{J_t}
$$
The classification of open coverings of $Y$ given above shows that
this rule transforms open coverings into open coverings. Thus we obtain
an abelian sheaf $\mathcal{G}$ on $Y$ by setting
$\mathcal{G}(Y) = \mathcal{F}(U)$ and for
$V = \bigcup\nolimits_{t \in T} V_{J_t}$ setting
$$
\mathcal{G}(V) = \mathcal{F}\left(\bigcup\nolimits_{t \in T} U_{J_t}\right)
$$
and using the restriction maps of $\mathcal{F}$.

\medskip\noindent
With these preliminaries out of the way we can prove our lemma as follows.
We have an open covering
$\mathcal{V} : Y = \bigcup_{i \in I} V_{\{i\}}$ of $Y$.
By construction we have an equality
$$
\check{C}^\bullet(\mathcal{V}, \mathcal{G}) =
\check{C}^\bullet(\mathcal{U}, \mathcal{F})
$$
of {\v C}ech complexes. Since the sheaf $\mathcal{G}$ is flasque on $Y$
(by our assumption on $\mathcal{F}$ in the statement of the lemma)
the vanishing follows from
Lemma \href{https://stacks.math.columbia.edu/tag/09SZ}{lemma flasque acyclic cech (Tag 09SZ)}.
\end{proof}


\begin{lemma}
\label{lemma-cech-vanish-basis}
(Variant of Lemma \href{https://stacks.math.columbia.edu/tag/01EV}{lemma cech vanish (Tag 01EV)}.)
Let $X$ be a ringed space.
Let $\mathcal{B}$ be a basis for the topology on $X$.
Let $\mathcal{F}$ be an $\mathcal{O}_X$-module.
Assume there exists a set of open coverings $\text{Cov}$
with the following properties:
\begin{enumerate}
\item For every $\mathcal{U} \in \text{Cov}$
with $\mathcal{U} : U = \bigcup_{i \in I} U_i$ we have
$U, U_i \in \mathcal{B}$ and every $U_{i_0 \ldots i_p} \in \mathcal{B}$.
\item For every $U \in \mathcal{B}$ the open coverings of $U$
occurring in $\text{Cov}$ is a cofinal system of open coverings
of $U$.
\item For every $\mathcal{U} \in \text{Cov}$ we have
$\check{H}^p(\mathcal{U}, \mathcal{F}) = 0$ for all $p > 0$.
\end{enumerate}
Then $H^p(U, \mathcal{F}) = 0$ for all $p > 0$ and any $U \in \mathcal{B}$.
\end{lemma}

\begin{proof}
Let $\mathcal{F}$ and $\text{Cov}$ be as in the lemma.
We will indicate this by saying ``$\mathcal{F}$ has vanishing higher
{\v C}ech cohomology for any $\mathcal{U} \in \text{Cov}$''.
Choose an embedding $\mathcal{F} \to \mathcal{I}$ into an
injective $\mathcal{O}_X$-module.
By Lemma \href{https://stacks.math.columbia.edu/tag/01EP}{lemma injective trivial cech (Tag 01EP)} $\mathcal{I}$
has vanishing higher {\v C}ech cohomology for any $\mathcal{U} \in \text{Cov}$.
Let $\mathcal{Q} = \mathcal{I}/\mathcal{F}$
so that we have a short exact sequence
$$
0 \to \mathcal{F} \to \mathcal{I} \to \mathcal{Q} \to 0.
$$
By Lemma \href{https://stacks.math.columbia.edu/tag/01EU}{lemma ses cech h1 (Tag 01EU)} and our assumption (2)
this sequence gives rise to an exact sequence
$$
0 \to \mathcal{F}(U) \to \mathcal{I}(U) \to \mathcal{Q}(U) \to 0.
$$
for every $U \in \mathcal{B}$. Hence for any $\mathcal{U} \in \text{Cov}$
we get a short exact sequence of {\v C}ech complexes
$$
0 \to
\check{\mathcal{C}}^\bullet(\mathcal{U}, \mathcal{F}) \to
\check{\mathcal{C}}^\bullet(\mathcal{U}, \mathcal{I}) \to
\check{\mathcal{C}}^\bullet(\mathcal{U}, \mathcal{Q}) \to 0
$$
since each term in the {\v C}ech complex is made up out of a product of
values over elements of $\mathcal{B}$ by assumption (1).
In particular we have a long exact sequence of {\v C}ech cohomology
groups for any open covering $\mathcal{U} \in \text{Cov}$.
This implies that $\mathcal{Q}$ is also an $\mathcal{O}_X$-module
with vanishing higher {\v C}ech cohomology for all
$\mathcal{U} \in \text{Cov}$.

\medskip\noindent
Next, we look at the long exact cohomology sequence
$$
\xymatrix{
0 \ar[r] &
H^0(U, \mathcal{F}) \ar[r] &
H^0(U, \mathcal{I}) \ar[r] &
H^0(U, \mathcal{Q}) \ar[lld] \\
&
H^1(U, \mathcal{F}) \ar[r] &
H^1(U, \mathcal{I}) \ar[r] &
H^1(U, \mathcal{Q}) \ar[lld] \\
&
\ldots & \ldots & \ldots \\
}
$$
for any $U \in \mathcal{B}$. Since $\mathcal{I}$ is injective we
have $H^n(U, \mathcal{I}) = 0$ for $n > 0$ (see
Derived Categories, Lemma \href{https://stacks.math.columbia.edu/tag/015B}{lemma higher derived functors (Tag 015B)}).
By the above we see that $H^0(U, \mathcal{I}) \to H^0(U, \mathcal{Q})$
is surjective and hence $H^1(U, \mathcal{F}) = 0$.
Since $\mathcal{F}$ was an arbitrary $\mathcal{O}_X$-module with
vanishing higher {\v C}ech cohomology for all $\mathcal{U} \in \text{Cov}$
we conclude that also $H^1(U, \mathcal{Q}) = 0$ since $\mathcal{Q}$ is
another of these sheaves (see above). By the long exact sequence this in
turn implies that $H^2(U, \mathcal{F}) = 0$. And so on and so forth.
\end{proof}


\begin{proposition}
\label{proposition-cohomological-dimension-spectral}
\begin{reference}
Part (1) is the main theorem of \href{https://stacks.math.columbia.edu/bibliography}{Bibliography: Scheiderer}.
\end{reference}
Let $X$ be a spectral space of Krull dimension $d$.
Let $\mathcal{F}$ be an abelian sheaf on $X$.
\begin{enumerate}
\item $H^q(X, \mathcal{F}) = 0$ for $q > d$,
\item $H^d(X, \mathcal{F}) \to H^d(U, \mathcal{F})$ is surjective
for every quasi-compact open $U \subset X$,
\item $H^q_Z(X, \mathcal{F}) = 0$ for $q > d$ and any constructible
closed subset $Z \subset X$.
\end{enumerate}
\end{proposition}

\begin{proof}
We prove this result by induction on $d$.

\medskip\noindent
If $d = 0$, then $X$ is a profinite space, see
Topology, Lemma \href{https://stacks.math.columbia.edu/tag/0905}{lemma characterize profinite spectral (Tag 0905)}.
Thus (1) holds by Lemma \href{https://stacks.math.columbia.edu/tag/0A3F}{lemma vanishing for profinite (Tag 0A3F)}.
If $U \subset X$ is quasi-compact open, then $U$ is
also closed as a quasi-compact subset of a Hausdorff space.
Hence $X = U \amalg (X \setminus U)$ as a topological space
and we see that (2) holds. Given $Z$ as in (3) we consider the
long exact sequence
$$
H^{q - 1}(X, \mathcal{F}) \to
H^{q - 1}(X \setminus Z, \mathcal{F}) \to
H^q_Z(X, \mathcal{F}) \to H^q(X, \mathcal{F})
$$
Since $X$ and $U = X \setminus Z$ are profinite (namely $U$ is quasi-compact
because $Z$ is constructible) and since
we have (2) and (1) we obtain the desired vanishing of the
cohomology groups with support in $Z$.

\medskip\noindent
Induction step. Assume $d \geq 1$ and assume
the proposition is valid for all spectral
spaces of dimension $< d$. We first prove part (2) for $X$.
Let $U$ be a quasi-compact open. Let $\xi \in H^d(U, \mathcal{F})$.
Set $Z = X \setminus U$. Let $W \subset X$ be the set of points
specializing to $Z$. By
Lemma \href{https://stacks.math.columbia.edu/tag/0A3D}{lemma cohomology of neighbourhoods of closed (Tag 0A3D)} we have
$$
H^d(W \setminus Z, \mathcal{F}|_{W \setminus Z}) =
\colim_{Z \subset V} H^d(V \setminus Z, \mathcal{F})
$$
where the colimit is over the quasi-compact open neighbourhoods $V$
of $Z$ in $X$.
By Topology, Lemma \href{https://stacks.math.columbia.edu/tag/0A31}{lemma make spectral space (Tag 0A31)} we see that
$W \setminus Z$ is a spectral space.
Since every point of $W$ specializes to a point of $Z$, we see that
$W \setminus Z$ is a spectral space of Krull dimension $< d$.
By induction hypothesis we see that the image of $\xi$ in
$H^d(W \setminus Z, \mathcal{F}|_{W \setminus Z})$ is zero.
By the displayed formula, there exists a $Z \subset V \subset X$
quasi-compact open such that $\xi|_{V \setminus Z} = 0$.
Since $V \setminus Z = V \cap U$ we conclude by the Mayer-Vietoris
(Lemma \href{https://stacks.math.columbia.edu/tag/01EB}{lemma mayer vietoris (Tag 01EB)}) for the covering $X = U \cup V$
that there exists a $\tilde \xi \in H^d(X, \mathcal{F})$ which restricts
to $\xi$ on $U$ and to zero on $V$. In other words, part (2) is true.

\medskip\noindent
Proof of part (1) assuming (2). Choose an injective resolution
$\mathcal{F} \to \mathcal{I}^\bullet$. Set
$$
\mathcal{G} = \Im(\mathcal{I}^{d - 1} \to \mathcal{I}^d) =
\Ker(\mathcal{I}^d \to \mathcal{I}^{d + 1})
$$
For $U \subset X$ quasi-compact open we have a map of exact sequences
as follows
$$
\xymatrix{
\mathcal{I}^{d - 1}(X) \ar[r] \ar[d] &
\mathcal{G}(X) \ar[r] \ar[d] &
H^d(X, \mathcal{F}) \ar[d] \ar[r] & 0 \\
\mathcal{I}^{d - 1}(U) \ar[r] &
\mathcal{G}(U) \ar[r] &
H^d(U, \mathcal{F}) \ar[r] & 0
}
$$
The sheaf $\mathcal{I}^{d - 1}$ is flasque by
Lemma \href{https://stacks.math.columbia.edu/tag/09SX}{lemma injective flasque (Tag 09SX)} and the fact that $d \geq 1$.
By part (2) we see that the right vertical arrow is surjective.
We conclude by a diagram chase that the map
$\mathcal{G}(X) \to \mathcal{G}(U)$ is surjective.
By Lemma \ref{lemma-vanishing-ravi} we conclude that
$\check{H}^q(\mathcal{U}, \mathcal{G}) = 0$ for $q > 0$ and
any finite covering $\mathcal{U} : U = U_1 \cup \ldots \cup U_n$
of a quasi-compact open by quasi-compact opens. Applying
Lemma \ref{lemma-cech-vanish-basis} we find that $H^q(U, \mathcal{G}) = 0$
for all $q > 0$ and all quasi-compact opens $U$ of $X$.
By Leray's acyclicity lemma
(Derived Categories, Lemma \href{https://stacks.math.columbia.edu/tag/015E}{lemma leray acyclicity (Tag 015E)})
we conclude that
$$
H^q(X, \mathcal{F}) =
H^q\left(
\Gamma(X, \mathcal{I}^0) \to \ldots \to
\Gamma(X, \mathcal{I}^{d - 1}) \to \Gamma(X, \mathcal{G})
\right)
$$
In particular the cohomology group vanishes if $q > d$.

\medskip\noindent
Proof of (3).  Given $Z$ as in (3) we consider the long exact sequence
$$
H^{q - 1}(X, \mathcal{F}) \to
H^{q - 1}(X \setminus Z, \mathcal{F}) \to
H^q_Z(X, \mathcal{F}) \to H^q(X, \mathcal{F})
$$
Since $X$ and $U = X \setminus Z$ are spectral spaces
(Topology, Lemma \href{https://stacks.math.columbia.edu/tag/0902}{lemma spectral sub (Tag 0902)})
of dimension $\leq d$
and since we have (2) and (1) we obtain the desired vanishing.
\end{proof}
\end{document}
